\chapter{Stool Elastase}
\section*{What Is This Test?}
Fecal elastase-1 (FE-1) is a pancreatic enzyme test that measures the concentration of elastase-1, a proteolytic enzyme produced by pancreatic acinar cells, in a single stool sample. Elastase-1 is a pancreas-specific enzyme that is not degraded during intestinal transit, and its concentration in stool is directly proportional to pancreatic exocrine secretory capacity. Fecal elastase-1 is the preferred noninvasive screening test for pancreatic exocrine insufficiency (PEI), a condition in which the pancreas fails to secrete adequate amounts of digestive enzymes (amylase, lipase, proteases including elastase), resulting in maldigestion of fats, proteins, and carbohydrates. PEI presents clinically with steatorrhea (bulky, greasy, foul-smelling stools that float and are difficult to flush), weight loss despite adequate caloric intake, fat-soluble vitamin deficiencies (A, D, E, K), and malnutrition. Causes include: chronic pancreatitis (most common cause in adults, typically from alcohol), cystic fibrosis (most common cause in children, CFTR mutations causing thick pancreatic secretions that obstruct ducts, leading to progressive acinar destruction and fibrosis), pancreatic ductal obstruction (pancreatic cancer, particularly adenocarcinoma of the pancreatic head), pancreatic resection (Whipple procedure, distal pancreatectomy), Shwachman-Diamond syndrome, and autoimmune pancreatitis. Fecal elastase-1 is stable in stool, unaffected by concurrent pancreatic enzyme replacement therapy [PERT] (the assay uses monoclonal antibodies that detect human elastase-1 specifically, not porcine elastase in PERT), and can be measured from a single random stool sample.
\section*{Alternative Names}
Fecal Elastase-1; FE-1; Pancreatic Elastase; Stool Elastase; Fecal Pancreatic Elastase
\section*{Principle}
Fecal elastase-1 is measured by ELISA (enzyme-linked immunosorbent assay) using two monoclonal antibodies directed against human pancreatic elastase-1. The specimen is a random stool sample (approximately 1--5 g, solid or liquid). Elastase-1 is extracted from stool and quantified by the sandwich ELISA; results are reported in mcg elastase/g stool. Porcine pancreatic elastase (present in PERT) does not cross-react with the human-specific monoclonal antibodies used in the assay, so PERT does NOT need to be discontinued before testing.
\section*{History}
Pancreatic exocrine insufficiency has been recognized since the nineteenth century, but reliable noninvasive tests were lacking for most of the twentieth century; the secretin--cholecystokinin stimulation test with duodenal intubation served as the reference standard but was invasive and rarely performed. Fecal elastase-1 was described in the mid-1990s, when Loser and colleagues (1996) and Dominguez-Munoz and colleagues (1995) independently showed that stool elastase-1 concentrations reflected pancreatic exocrine output and remained stable during intestinal transit. Because elastase-1 is not degraded in the gut and is unaffected by pancreatic enzyme replacement therapy, it quickly became the most widely used noninvasive test for pancreatic insufficiency. Subsequent studies established the widely used cutoffs of greater than 200 mcg/g (normal), 100--200 mcg/g (moderate insufficiency), and less than 100 mcg/g (severe insufficiency). Fecal elastase-1 is now recommended by gastroenterology guidelines for the initial diagnosis of PEI in chronic pancreatitis, cystic fibrosis, and after pancreatic resection.
\section*{Why This Test Is Ordered}
\begin{itemize}
  \item Evaluation of suspected pancreatic exocrine insufficiency in patients with steatorrhea, chronic diarrhea, weight loss, or malnutrition
  \item Diagnosis and monitoring of PEI in chronic pancreatitis, including alcoholic and hereditary forms
  \item Assessment of pancreatic function in cystic fibrosis, helping guide pancreatic enzyme replacement therapy (PERT)
  \item Screening in patients with cystic fibrosis-related diabetes or unexplained recurrent abdominal pain
  \item Evaluation after pancreatic surgery (Whipple procedure, distal pancreatectomy) to detect and manage postoperative PEI
  \item Assessment of suspected pancreatic ductal obstruction from pancreatic cancer or autoimmune pancreatitis
  \item Investigation of fat-soluble vitamin deficiencies or failure to thrive in children
\end{itemize}
\section*{Primary vs Secondary Test}
Fecal elastase-1 is the primary noninvasive screening test for pancreatic exocrine insufficiency. The 72-hour fecal fat quantification remains the gold standard confirmatory test for steatorrhea, and direct (secretin-stimulated) pancreatic function testing is reserved for specialized centers when results are equivocal.
\section*{How This Test Is Performed}
The patient provides a single random stool sample (approximately 1--5 g) collected directly into a clean, dry container; formed or semi-formed stool is preferred because watery stool can dilute elastase and produce falsely low results. In the laboratory, elastase-1 is extracted from the stool and measured by a sandwich ELISA using two monoclonal antibodies specific for human pancreatic elastase-1, and the result is reported in mcg elastase/g stool. Patients do not need to stop pancreatic enzyme replacement therapy, because the assay detects only human elastase, not the porcine enzyme contained in PERT.
\section*{What Preparation Is Needed}
No fasting, dietary restrictions, or discontinuation of pancreatic enzyme replacement therapy (PERT) are required, since the assay is specific for human elastase-1. The patient should collect a formed or semi-formed stool sample, avoiding contamination with water or toilet bowl contents, which can dilute the specimen. If the stool is watery or the result is unexpectedly low, the test should be repeated on a formed specimen.
\section*{Contraindications and Precautions}
There are no contraindications to stool collection. Precautions include using a formed specimen, since watery or dilute stool is the most common cause of a falsely low elastase result. The sample should not be contaminated with urine, water, or toilet bowl cleaner, and it should be transported promptly; elastase-1 is stable for at least several days at room temperature, but refrigeration is recommended if transport is delayed. Borderline reductions (100--200 mcg/g) require clinical correlation and may warrant confirmatory testing, because asymptomatic individuals can have borderline low values.
\section*{Risks}
No significant risks; the test is entirely noninvasive and requires only collection of a stool sample.
\section*{What Happens After the Test}
Results are usually available within a few days. A fecal elastase-1 below 200 mcg/g, especially below 100 mcg/g, confirms pancreatic exocrine insufficiency and supports starting PERT with each meal and snack, along with fat-soluble vitamin replacement and dietary counseling. If results are borderline or inconsistent with the clinical picture, the test is repeated on a formed specimen or followed by 72-hour fecal fat quantification and pancreatic imaging (CT, MRCP, or endoscopic ultrasound).
\section*{Understanding Results}
\subsection*{Severe Pancreatic Insufficiency (<100 mcg/g)}
Severely reduced or absent pancreatic exocrine function. Common in: cystic fibrosis, advanced chronic pancreatitis, pancreatic cancer with ductal obstruction, and after extensive pancreatic resection. Fecal elastase <100 mcg/g in a patient with steatorrhea and weight loss is diagnostic of severe PEI. Pancreatic enzyme replacement therapy (PERT) with enteric-coated microsphere/microtablet preparations (Creon, Zenpep, Pancreaze) with each meal and snack is indicated.
\subsection*{Moderate Pancreatic Insufficiency (100--200 mcg/g)}
Mild to moderate reduction in pancreatic function. May represent early/mild chronic pancreatitis, partial ductal obstruction, or early CF. Clinical correlation is required. PERT may be considered in symptomatic patients.
\subsection*{Normal Exocrine Function (>200 mcg/g)}
Normal pancreatic exocrine function. PEI is unlikely.
\subsection*{False-Positive Decrease}
Watery/dilute stool (diarrhea, laxatives, liquid stool) --- elastase concentration may be artifactually decreased by dilution. A formed/semi-formed stool specimen is preferred. This is the most common cause of a falsely low fecal elastase result. If the fecal elastase is low (<200 mcg/g) and the stool is watery, the test should be repeated on a formed specimen.
\section*{Normal Results}
Normal: >200 mcg/g stool (pancreatic sufficiency). Moderate PEI: 100--200 mcg/g. Severe PEI: <100 mcg/g. Fecal elastase-1 >500 mcg/g is typical of healthy individuals. The test is NOT affected by PERT (pancreatic enzyme replacement therapy) --- patients do not need to hold enzymes before testing.
\section*{Follow-up Tests}
72-hour fecal fat quantification (the gold standard for steatorrhea; >7 g fat/24 hours on a 100 g fat/day diet is diagnostic of fat malabsorption). Fecal fat Sudan stain qualitative assessment. Serum trypsinogen (low in CF). Pancreatic imaging: CT abdomen with IV contrast, MRI/MRCP (magnetic resonance cholangiopancreatography), endoscopic ultrasound (EUS). Secretin-enhanced MRCP or endoscopic secretin pancreatic function testing (research/respecialized centers). Cystic fibrosis: sweat chloride test, CFTR genetic testing.
\section*{References}
\begin{enumerate}
  \item Loser C, Mollgaard A, Folsch UR. Faecal elastase 1: a novel, highly sensitive, and specific tubeless pancreatic function test. Gut. 1996;39(4):580--586.
  \item Dominguez-Munoz JE, Hieronymus C, Sauerbruch T, Malfertheiner P. Fecal elastase test: evaluation of a new noninvasive pancreatic function test. Am J Gastroenterol. 1995;90(10):1834--1837.
  \item Lindkvist B. Diagnosis and treatment of pancreatic exocrine insufficiency. World J Gastroenterol. 2013;19(42):7258--7266.
  \item Struyvenberg MR, Martin CR, Freedman SD. Practical guide to exocrine pancreatic insufficiency: breaking the myths. BMC Medicine. 2017;15(1):29.
  \item Layer P, Keller J. Lipase supplementation therapy: standards, alternatives, and perspectives. Pancreas. 2003;26(1):1--7.
\end{enumerate}
\clearpage
